
et on
définit un homomorphisme :
\[
  u:K\longrightarrow K_1\times K'
\]
en posant :
\[
  u(y)=\bigl(g^*(y).\lambda_{-1}(F),-i_*(y)\bigr),
\]
avec \(F=N-L\).

\textbf{Proposition 4.7.5.} \uline{Les opérations définies par} (4.7.1) \uline{et} (4.7.2) \uline{munissent} \(K_1\times K'\) \uline{d'une structure de} \(\lambda\)-\uline{anneau}, \uline{et} \(u\) \uline{est un homomorphisme injectif qui identifie} \(K\), \uline{muni des opérations définies par} (4.7.3) \uline{et} (4.7.4) \uline{à un} \(\lambda\)-\uline{idéal de} \(K_1\times K'\).

Il est immédiat de vérifier que \(K_1\times K'\) est un \(\lambda\)-anneau : on se réfèrera par exemple à V 5.5, où une construction tout à fait analogue est faite. La démonstration de la seconde partie de la proposition se fait également sans difficultés en utilisant les relations précédentes, la relation \(i_*(y_1y_2\lambda_{-1}(N))=i_*(y_1).i_*(y_2)\) qui résulte immédiatement des hypothèses, la relation \(g_*(\lambda_{-1}(F))=1\) (voir VI 5.9) et les relations de V 5.

On pose alors :
\[
  K'_1=(K_1\times K')/K
\]
et on note \(v\) l'homomorphisme canonique \(K_1\times K'\to K'_1\). On considère \(K'_1\) comme \(\lambda\)-anneau grâce à la structure quotient.

On définit ensuite des homomorphismes additifs :
\[
\begin{array}{llll}
j_*:K_1\to K'_1; & j^*:K'_1\to K_1; \\
f_*:K'_1\to K'; & f^*:K'\to K'_1 .
\end{array}
\]

L'homomorphisme \(j_*\) est le composé \(K_1\to K_1\times K'\to K'_1\). De même, l'homomorphisme \(f^*\) est le composé \(K'\to K_1\times K'\to K'_1\). On a alors :
\[
  v=j_*+f^*.
\]
